\section{Background}

\label{sec:background}

\subsection{ASO Experiential Priors (Recap)}

In ASO~\cite{hanzo2026aso}, each experiential prior $E_m$ consists of:
\begin{itemize}[leftmargin=1.5em]
    \item A 1-bit quantized expert factor $\widehat{\Delta}_m \in \{-1, +1\}^{|\mathcal{V}_m|}$
          with scale $\alpha_m > 0$.
    \item A quality score $q_m \in (0, 1)$ estimating pattern reliability.
    \item A context matching predicate $C_m(x, y_{<t}) \in \{0, 1\}$ specifying
          when the prior applies.
    \item Metadata: task embedding $\bm{e}_m \in \mathbb{R}^d$, timestamp,
          source attribution.
\end{itemize}

At decode time, priors contribute to the Product-of-Experts distribution:
\begin{equation}
    \pi_{\text{ASO}}(y_t \mid x, y_{<t}) \propto \pi_\theta(y_t \mid x, y_{<t})
    \prod_{m=1}^{M} \phi_m(y_t \mid x, y_{<t})^{\eta_m},
    \label{eq:poe-recap}
\end{equation}
where $\eta_m = \log(q_m / (1 - q_m))$ and $\phi_m$ is constructed from
$\widehat{\Delta}_m$.

\subsection{Byzantine Fault Tolerance}

A distributed protocol is \emph{$f$-Byzantine-fault-tolerant} if it maintains
correctness properties when up to $f$ of $n$ participants behave
arbitrarily~\cite{lamport1982byzantine}. The classical bound for
agreement-based protocols is $f < n/3$~\cite{pease1980reaching}.

\begin{definition}[Byzantine Node]
A node $i$ is Byzantine if it may deviate from the protocol specification in any
way, including:
\begin{enumerate}[leftmargin=1.5em]
    \item Submitting arbitrary (possibly adversarial) prior data.
    \item Withholding or selectively propagating messages.
    \item Colluding with other Byzantine nodes.
    \item Creating multiple identities (Sybil attack).
\end{enumerate}
\end{definition}

\subsection{Weighted Median}

The weighted median is a robust estimator that tolerates up to half the total
weight in outliers~\cite{minsker2015geometric}.

\begin{definition}[Weighted Median]
Given values $\{x_i\}_{i=1}^n$ with non-negative weights $\{w_i\}_{i=1}^n$
summing to $W = \sum_i w_i$, the weighted median is:
\begin{equation}
    \wmed(\{x_i, w_i\}) = \inf\left\{x : \sum_{i: x_i \leq x} w_i \geq W/2\right\}.
\end{equation}
\end{definition}

The weighted median has a breakdown point of $1/2$ in total weight, meaning an
adversary must control more than half the total weight to move the estimate
arbitrarily.

% ============================================================================
% 3. THREE-LAYER STORAGE ARCHITECTURE
% ============================================================================
